\section{Introduction}

There is a particularly sharp way for a local-global principle to fail.
Take a finite \'etale scheme \(X\) over \(\Q\).  It has no singularities,
no positive-dimensional families of points, and no geometry in which a
rational point can hide.  Nevertheless it may satisfy
\[
 X(\R)\ne\varnothing,\qquad
 X(\Q_p)\ne\varnothing\quad\text{for every prime }p,
 \qquad
 X(\Q)=\varnothing.
 \tag{1.1}\label{eq:hasse-condition}
\]
For a monic polynomial this says something completely concrete: at each
completion at least one factor has a root, although none has a rational
root.  Such polynomials are called intersective.

The classical example of Berend and Bilu is
\[
 f_{19}(X)=(X^3-19)(X^2+X+1).
 \tag{1.2}\label{eq:berend-bilu}
\]
The two factors divide the local work between them.  At primes
\(p\equiv2\pmod3\), cubing is invertible and the cubic supplies a root.  At
primes \(p\equiv1\pmod3\), the quadratic supplies a cube root of unity.
Only \(2,3,\) and \(19\) require separate attention.  Globally, neither
factor has a rational root.

The companion prescribed-fiber theorem turns every finite \'etale algebra
of rank at least three into a full fiber of a Jacobian-one map of affine
three-space \cite{vanRijnFibers}.  In particular, it turns
\eqref{eq:berend-bilu} into a degree-five Hasse failure.  At first sight
this seems to settle the Keller-map version of the problem.

It settles only the first quantifier order.  The prescribed-fiber
construction begins with a polynomial and builds a map around it.  If the
polynomial changes, the map generally changes as well.  The question that
remained was therefore:

\begin{quote}
\emph{Can one polynomial Keller map contain infinitely many Hasse-failing
complete fibers?}
\end{quote}

This is a different problem.  We no longer want to compile one arithmetic
algebra into one map.  We need an entire arithmetic family to fit inside the
target pencil of a single map.  That small change in quantifiers is where the
problem becomes interesting.

\subsection*{Finding the line}

The Berend--Bilu example already suggests the family
\[
 f_a(X)=(X^3-a)(X^2+X+1).
 \tag{1.3}\label{eq:intersective-family}
\]
For primes \(a=\ell\equiv1\pmod9\), the same local proof works uniformly.
But the expanded polynomial
\[
 f_a(X)=X^5+X^4+X^3-aX^2-aX-a
 \]
does not look like a one-coordinate motion in the usual linear inverse
pencil: three coefficients move at once.

The useful step is almost embarrassingly small.  Center the quadratic:
\[
 X=S-\frac12,\qquad
 X^2+X+1=S^2+\frac34.
 \tag{1.4}\label{eq:centered-quadratic-intro}
\]
Now the parameter enters through
\[
 -a\left(S^2+\frac34\right),
 \]
so only the quadratic and constant coefficients move.  Those are exactly
the two target coefficients in the quadratic-gauge inverse equation
\[
 E_{\Pi,B,C}(S)=G_\Pi(S)-\frac{g_1}{2}(BS^2+C).
 \tag{1.5}\label{eq:inverse-intro}
\]
Once \eqref{eq:centered-quadratic-intro} is made, the family
\eqref{eq:intersective-family} is no longer a search problem.  It is a
straight line in the target of one fixed map.  This was the point at which
the isolated example began to look like a family.

This gives the geometric part of the discovery.  The arithmetic part then
becomes unexpectedly generous.  Primality of \(a\) is not used in the local
argument.  What is really needed is
\[
 a\equiv1\pmod9,\qquad a\notin\Q^3,\qquad
 p\mid a\Longrightarrow p\equiv1\pmod3.
 \tag{1.6}\label{eq:parameter-conditions}
\]
Thus the prime progression sits inside a much larger multiplicatively
generated family.  Its congruence-and-support core is a semigroup; the
noncube condition removes the perfect cubes, so \(\mathcal A\) itself is
not literally closed under multiplication.  The prime family has order
\(B/\log B\) by height; the full sufficient family has order
\(B/\sqrt{\log B}\).

\subsection*{The result}

For \(a>1\), define
\[
 y_a=\left(-1,\frac{32a}{9},\frac{8a+1}{3}\right)
 \in\A^3(\Q)
 \tag{1.7}\label{eq:target-family}
\]
and let
\[
 \mathcal A=
 \left\{
 \begin{array}{@{}l@{}}
 a>1:\ a\equiv1\pmod9,\ a\notin\Q^3,\\
 \hfill p\mid a\Longrightarrow p\equiv1\pmod3
 \end{array}
 \right\}.
 \tag{1.8}\label{eq:parameter-family}
\]
Since \(a\) is an integer, the condition \(a\notin\Q^3\) simply excludes
perfect cubes.

\begin{theorem}[Fixed-map theorem]\label{thm:main}
There is an explicit polynomial map
\[
 \Phi:\A^3_{\Q}\longrightarrow\A^3_{\Q},
 \qquad \det D\Phi=1,\qquad \gdeg(\Phi)=5,
 \]
such that, for every \(a\in\mathcal A\),
\[
 \boxed{
 \Phi^{-1}(y_a)
 \simeq
 \Spec\Q[X]/\bigl((X^3-a)(X^2+X+1)\bigr).}
 \tag{1.9}\label{eq:main-fiber}
\]
This is a full reduced fiber.  It has a point over \(\R\) and over every
\(\Q_p\), but it has no \(\Q\)-point.
\end{theorem}

The targets have primitive projective coordinates
\[
 [1:y_a]=[9:-9:32a:24a+3],
 \qquad H(y_a)=32a.
 \tag{1.10}\label{eq:height-intro}
\]
They are therefore pairwise distinct.  The first parameter is \(a=19\), so
the classical example is the first point on the line rather than an
isolated construction.  The old example has not disappeared; it has found
its natural place.

\begin{theorem}[Quantitative theorem]\label{thm:count-intro}
Let
\[
 M_\Phi(B)=\#\{a\in\mathcal A:H(y_a)\le B\}.
 \]
Then
\[
 \boxed{
 M_\Phi(B)\sim
 \frac{G_3(1)}{96\sqrt{\pi}}\,
 \frac{B}{\sqrt{\log B}},}
 \tag{1.11}\label{eq:count-intro}
\]
where
\[
 G_3(1)=
 \left[
 \frac23 L(1,\chi_{-3})
 \prod_{p\equiv2\ (3)}(1-p^{-2})
 \right]^{1/2}>0.
 \tag{1.12}\label{eq:g3-intro}
\]
Consequently the total number of Hasse-failing rational targets of
\(\Phi\) of height at most \(B\) is at least of order
\(B/\sqrt{\log B}\).
\end{theorem}

The theorem counts the displayed sufficient family exactly; it does not
claim an asymptotic for all Hasse-failing targets of \(\Phi\).  This
distinction matters.  The bounded experiment later in the paper checks the
family and the numerical convergence, but the asymptotic comes from the
Euler product and Selberg--Delange.

The rest of the paper follows the route by which the construction becomes
visible.  Section~\ref{sec:fixed-map} starts from the scalar inverse
equation, writes down the fixed map, and proves that no boundary point is
missing from the fiber.  Section~\ref{sec:local-global} lets the two factors
in \eqref{eq:intersective-family} divide the completions between them.
Section~\ref{sec:counting} first counts the prime progression and then the
larger multiplicative family.  Section~\ref{sec:minimality} returns to the
classical degree obstruction and shows that five is not only attained, but
forced.
